\section{Simulation setup}
For the present study, we use configurations generated with   2+1 flavor clover fermions and tree-level Symanzik gauge action configuration by the CLS collaboration using periodic boundary conditions~\cite{Bruno:2014jqa}. The detailed parameters are listed in Table~\ref{table:cls}. Note that $m_\pi=547$ MeV instead of 333 MeV is used for valence quarks in order to have a better signal. Physically,  the soft function
becomes  independent of the meson mass for large boost factors $\gamma$.

To calculate the form factor in Eq.~(\ref{eq:FF}), we generate the wall source propagator,
\begin{align}
 S_w(x,t,t';\vec{p})=\sum_{\vec{y}} S(t,\vec{x};t',\vec{y})e^{i\vec{p}\cdot(\vec{y}-\vec{x})},
\end{align}
on the Coulomb gauge fixed configurations at $t'=0$ and $t_{\rm sep}$ for both the initial and final meson states.  $S$ is the quark propagator from $(t',\vec{y})$ to $(t,\vec{x})$. Then we can construct the three point function (3pt)
corresponding to the form factor in Eq.~(\ref{eq:FF}),
\begin{align}
&C_3(b_\perp,P^z;p^z,t_{\rm sep},t)\\
&=\frac{1}{L^3} \sum_x\textrm{Tr} \langle  S_{w}^\dagger(\vec{x}+\vec{b},t,0;-\vec{p})\gamma_5\Gamma S_{w}(\vec{x}+\vec{b},t,t_{\rm sep};\vec{p}) \nonumber\\
&\;\;\quad\quad \quad \times S_{w}^\dagger(\vec{x},t,t_{\rm sep};-\vec{P}+\vec{p}) \gamma_{5}\Gamma
 S_{w}(\vec{x},t,0;\vec{P}-\vec{p})\rangle\nonumber.
\end{align}
The quark momentum $\vec{p}=(\vec 0_{\perp},p^z)$, and
the relation $\gamma_5 S^{\dagger}(x,y)\gamma_5=S(y,x)$ have been applied for the anti-quark propagator. We have tested several choices of $\Gamma$, and will use the unity Dirac matrix  $\Gamma=I$ as it has the best signal and describes the leading twist light-cone contribution in the large $P^z$ limit. Notice that the $\Gamma=\gamma_4$ case  is   subleading in the large $P^z$ limit and is less suitable,   although  the excited state contamination might be smaller.

By generating the wall source propagators at all the 48 time slices with quark momentum ${p}^z=(-2,-1,0,1,2) \times 2\pi/(La)$, we can maximize the statistics of the 3pt function with all the meson momenta $P_z$ from 0 to $8\pi/(La)$ ($\sim 2.1$ GeV) with arbitrary $t$ and $t_{\rm sep}$.  $C_3(b_\perp,P^z,t_{\rm sep},t)$ is related to the bare $F(b_\perp,P^z)$ using standard parameterization of 3pt with one excited {state},
\begin{align}
& C_3(b_\perp,P^z;p^z,t_{\rm sep},t)=\frac{A_w(p_z)^2}{(2E)^2}e^{-Et_{\rm sep}}\big[F(b_\perp,P^z)\nonumber\\
&\quad\quad+c_1(e^{-\Delta E t}+e^{-\Delta E (t_{\rm sep}-t)})+c_2e^{-\Delta E t_{\rm sep}}\big].
 \label{eq:3pt}
\end{align}
$A_w$ is the matrix element of the Coulomb gauge fixed wall (CFW) source pion interpolation field, $E=\sqrt{m_{\pi}^2+{P}^{z2}}$ is the pion energy, $\Delta E$ is the mass gap between pion and its first excited state, $c_{1,2} $ are parameters for the excited state contamination. Note that the $p_z$ dependence factor $A_w^2$ will cancel{.}

The same wall source propagators can be used to calculate the two-point function related to the bare quasi-TMDWF,
\begin{align}
&C_2(b_\perp,P^z;p_z,\ell,t)=\frac{1}{L^3\sqrt{Z_E(2\ell,b_\perp)}}\sum_x \textrm{Tr}e^{i\vec{P}\cdot \vec{x}}\nonumber\\
&\times\langle  S_w^\dagger(\vec{x}+\vec{b},t,0;-\vec{p}){\cal W}(\vec b,\ell) \gamma_5\Gamma_\Phi  S_w(\vec{x},t,0;P^z-\vec{p}) \rangle\nonumber\\
& =\frac{A_w(p_z)A_p}{2E}e^{-Et}\phi_\ell(0,b_\perp,P^z,\ell)(1+c_0 e^{-\Delta E t}),\label{eq:2pt}
\end{align}
where again we parameterize the mixing with one excited state. 
$A_p$ is the matrix element of the point sink pion interpolation field. It will be  removed when we normalize $\phi_\ell(0,b_\perp,P^z,\ell)$ with $\phi_\ell(0,0,P^z,0)$. We choose $\Gamma_\Phi=\gamma^t\gamma_5$
to define the wave function amplitude in Eq.~(\ref{eq:phi_def}). Based on the quasi-TMDPDF study in Ref.{~\cite{Shanahan:2019zcq,Shanahan:2020zxr}} with a similar staple-shaped gauge link operator, the mixing effect {could be sizable  when summing various contributions. In the supplemental material, we report a similar simulation but using the A654 ensemble. We find that the mixing effects can reach   order $5\%$ for the transverse separation $b_\perp\sim 0.6{\rm fm}$. These effects will be included in the following analysis as  one of the systematic uncertainties, while a comprehensive study on the mixing effects will be conducted in the future. } 


\begin{figure}[!th]
\includegraphics[width=0.45\textwidth]{images/02_tmdwf_mod_with_L_P3_b3_t6.pdf}
\caption{ Results for the $\ell$ dependence of the  quasi-TMDWF with $z=0$, and also the square root of the Wilson loop which is  used for the subtraction, taking the $\{P^z, b_{\perp}, t\}=\{6\pi/L, 3a, 6a\}$ case as a example. All the results are normalized with their values at $\ell=0$.}\label{fig:L_dependence}
\end{figure}

The dispersion relation of the pion state, statistical checks for the  measurement histogram, and information on the autocorrelation between configurations can be found in the supplemental materials~\cite{supplemental}.
